\section{Bridges in a Graph}
\label{sec:graph-bridges}

\problemstatement{Given an undirected, connected graph, find every
\textbf{bridge} -- an edge whose removal would increase the number of
connected components (disconnect the graph).}

\begin{patternbox}
\emph{``Find edges whose removal disconnects the graph''} is solved by
a single DFS tracking $\textit{tin}$ and $\textit{low}$
(Section~\ref{sec:graph-others-intro}): a tree edge $u \to v$ (where
$v$ is $u$'s child in the DFS tree) is a bridge exactly when
$\textit{low}[v] > \textit{tin}[u]$ -- meaning $v$'s entire subtree
has \textbf{no} back edge reaching $u$ or any of $u$'s ancestors, so
the only way to reach $v$'s subtree from the rest of the graph is
through this one edge.
\end{patternbox}

\subsection*{Deriving the \textit{low}-value update rule}

While DFS explores from node $u$, for each neighbor $v$:
\begin{itemize}
  \item If $v$ is $u$'s \textbf{parent} (the edge just arrived
        through), skip it entirely -- this is not a genuine
        alternative path.
  \item If $v$ is \textbf{unvisited}, it becomes a tree-edge child:
        recurse, then update $\textit{low}[u] \gets
        \min(\textit{low}[u], \textit{low}[v])$ (whatever $v$'s
        subtree can reach back to, $u$ can also reach back to, through
        $v$). Check the bridge condition here:
        $\textit{low}[v] > \textit{tin}[u]$.
  \item If $v$ is \textbf{already visited} (and not the parent), this
        is a back edge: update $\textit{low}[u] \gets
        \min(\textit{low}[u], \textit{tin}[v])$ -- using $v$'s
        \textbf{discovery time} directly (not $\textit{low}[v]$, since
        crossing a back edge to $v$ only guarantees reaching $v$
        itself, not necessarily anything $v$'s own low value might
        represent through edges not yet explored from $u$'s
        perspective).
\end{itemize}

\subsection*{C++ implementation}

\begin{lstlisting}
#include <bits/stdc++.h>
using namespace std;

int timer = 0;

void dfs(int node, int parent, vector<vector<int>>& adj, vector<bool>& visited,
         vector<int>& tin, vector<int>& low, vector<vector<int>>& bridges) {
    visited[node] = true;
    tin[node] = low[node] = timer++;

    for (int neighbor : adj[node]) {
        if (neighbor == parent) continue;

        if (visited[neighbor]) {
            low[node] = min(low[node], tin[neighbor]); // back edge
        } else {
            dfs(neighbor, node, adj, visited, tin, low, bridges);
            low[node] = min(low[node], low[neighbor]);
            if (low[neighbor] > tin[node]) {
                bridges.push_back({node, neighbor});
            }
        }
    }
}

vector<vector<int>> findBridges(int n, vector<vector<int>>& adj) {
    vector<bool> visited(n, false);
    vector<int> tin(n), low(n);
    vector<vector<int>> bridges;
    timer = 0;

    for (int i = 0; i < n; i++) {
        if (!visited[i]) dfs(i, -1, adj, visited, tin, low, bridges);
    }
    return bridges;
}

int main() {
    int n = 4;
    vector<vector<int>> adj(n);
    auto addEdge = [&](int u, int v) { adj[u].push_back(v); adj[v].push_back(u); };
    addEdge(0, 1); addEdge(1, 2); addEdge(2, 0); addEdge(1, 3);

    for (auto& b : findBridges(n, adj)) cout << b[0] << "-" << b[1] << " ";
    cout << endl; // 1-3
    return 0;
}
\end{lstlisting}

\subsection*{Dry run}

\dryrun{Take edges $0{-}1,1{-}2,2{-}0,1{-}3$ (a triangle $0,1,2$ with
a pendant edge to $3$), DFS from $0$.}

$\textit{dfs}(0,-1)$: $\textit{tin}[0]=\textit{low}[0]=0$. Neighbor
$1$: unvisited, recurse. $\textit{dfs}(1,0)$:
$\textit{tin}[1]=\textit{low}[1]=1$. Neighbor $0$: is parent, skip.
Neighbor $2$: unvisited, recurse. $\textit{dfs}(2,1)$:
$\textit{tin}[2]=\textit{low}[2]=2$. Neighbor $1$: is parent, skip.
Neighbor $0$: already visited, not parent -- back edge,
$\textit{low}[2]=\min(2,\textit{tin}[0]{=}0)=0$. Return to $1$:
$\textit{low}[1]=\min(1,\textit{low}[2]{=}0)=0$; check bridge:
$\textit{low}[2]{=}0 > \textit{tin}[1]{=}1$? No -- not a bridge
(correct: edge $1{-}2$ sits on the triangle's cycle). Back at $1$,
next neighbor $3$: unvisited, recurse. $\textit{dfs}(3,1)$:
$\textit{tin}[3]=\textit{low}[3]=3$; no unvisited or back-edge
neighbors. Return to $1$: $\textit{low}[1]=\min(0,\textit{low}[3]{=}3)=0$;
check bridge: $\textit{low}[3]{=}3 > \textit{tin}[1]{=}1$? \textbf{Yes}
-- edge $1{-}3$ is a bridge (correct: it's the only connection to the
otherwise-isolated node $3$). Return to $0$: $\textit{low}[0]=
\min(0,\textit{low}[1]{=}0)=0$; check bridge:
$\textit{low}[1]{=}0 > \textit{tin}[0]{=}0$? No -- not a bridge.

\begin{complexitybox}
\textbf{Time Complexity.} $O(V+E)$ -- a single DFS pass.
\textbf{Space Complexity.} $O(V)$ for \textit{tin}, \textit{low},
\textit{visited}, and the recursion stack.
\end{complexitybox}

\begin{takeawaybox}
The distinction between updating $\textit{low}$ with a child's
$\textit{low}$ value (tree edges) versus a neighbor's $\textit{tin}$
value (back edges) is the single most important detail to get right in
this algorithm -- using $\textit{low}$ for a back edge would
incorrectly let information ``leak'' through a node that hasn't
finished exploring its own subtree yet.
\end{takeawaybox}
